\section{Classical support substitution before Gaussian integration}
\label{cdc-section}

This section places the classical precision and base factors directly
in the next background $w$, before applying the Gaussian expansion.
The composite ratio $R=\ell^s$ is fixed. Constants may depend on
$R$ and on the later auxiliary grouping, but are uniform in the
preceding homogeneous depth $M$ and compatible volume. The inputs
are the established classical strict families with arbitrary fixed
cardinality reserve and the actual supported stationary map. No
such arbitrary reserve is assumed for inherited quantum activities.

\subsection{A precision composition lemma}
\label{cdc-precision-section}

On one fixed cell graph let
$B(a)=\sum_Y B_Y(a|_Y)$ be a symmetric matrix family. Its matrix
rows and columns are supported in $Y$, and $B_Y(0)=0$. Each
$a$-cell has $D$ scalar coordinates with independent analytic
radius $b_{\rm cl}$. Write $\tau(Y)$ for graph Steiner-tree length,
allowing auxiliary vertices. Suppose
\begin{equation}
 \sup_i\sum_{Y\ni i}\sum_{j\in Y}
 e^{\sigma\tau(Y)+A|Y|}
       \|(B_Y)_{ij}\|_{b_{\rm cl}}\leq N_B.
 \label{cdc-precision-input}
\end{equation}
The norms inside this sum are block operator norms. Let the
stationary substitution have a normally convergent strict family
\begin{equation}
 a_i(w)=\sum_X u_{i;X}(w_X),\qquad u_{i;X}(0)=0,\qquad
 \sup_i\sum_Xe^{\kappa\tau(\{i\}\cup X)+a|X|}
                         \|u_{i;X}\|_r\leq s(r).
 \label{cdc-map-input}
\end{equation}
Here each output norm is the maximum over the entire cell fiber,
\emph{inside} the support sum.

\begin{lemma}[Composition in the matrix row norm]
\label{cdc-precision-lemma}
Choose $0<\kappa'<\min(\sigma,\kappa)$, $A'\geq0$ with
$a\geq A'$, and $\rho>0$. Put
\begin{equation}
 t=\frac{s(r)}{b_{\rm cl}}<1,\qquad
 \alpha_D(t)=D\log\frac1{1-t}.
 \label{cdc-alpha}
\end{equation}
If $A'+\alpha_D(t)+\rho\leq A$, then
$B(a(w))=\sum_Z B'_Z(w_Z)$ has symmetric strict pieces with
rows and columns in $Z$, and
\begin{equation}
 \sup_i\sum_{Z\ni i}\sum_{j\in Z}
 e^{\kappa'\tau(Z)+A'|Z|}\|(B'_Z)_{ij}\|_r
 \leq N_B\alpha_D(t)
                   \sup_{n\geq1}n e^{-\rho n}.
 \label{cdc-precision-output}
\end{equation}
In particular the bound is $O(r)$ if $s(r)=O(r)$ and the
displayed margins are fixed.
\end{lemma}
\begin{proof}
Expand a matrix entry of $B_Y$ on its scalar-coordinate polydisc.
Cauchy bounds its multiindex coefficient by
$\|(B_Y)_{ij}\|_{b_{\rm cl}}b_{\rm cl}^{-|\nu|}$.
The zero multiindex is absent. Substitute the strict family
\eqref{cdc-map-input} in every occurrence of every scalar input.
If an occurrence is rooted at $k\in Y$ and assigned support $X$,
declare the resulting carrier to be
$Z=Y\cup\bigcup X$. Steiner gluing gives
\begin{equation}
 \tau(Z)\leq\tau(Y)+\sum_{\rm occurrences}
                            \tau(\{k\}\cup X),\qquad
 |Z|\leq|Y|+\sum_{\rm occurrences}|X|.
 \label{cdc-precision-gluing}
\end{equation}
The scalar multiindex sum is
\begin{equation}
 (1-t)^{-D|Y|}-1
 \leq\alpha_D(t)|Y|e^{\alpha_D(t)|Y|}.
 \label{cdc-precision-generating}
\end{equation}
The $\rho$ reserve pays $|Y|$ and proves
\eqref{cdc-precision-output}. For a fixed nonzero matrix row
$i$, the original piece already had $i\in Y$, and its nonzero
columns remain in $Y$. Substitution does not create new matrix
rows or columns. Thus no path-anchor rerooting factor is needed.
Symmetry is preserved coefficientwise and after grouping equal
$Z$. Absolute convergence justifies all groupings.
\end{proof}

If only $\sum_YB_Y(0)=0$ is initially known, subtract the
individual constants. Their absolutely convergent sum is zero,
and the input norm increases by at most a factor two. The lemma
is a statement about the precision's weighted row norm. It does
not assert a norm of whole global matrices anchored at every
new support site.

\subsection{Application to the actual composite precision}
\label{cdc-actual-precision-section}

Group current-unit classical families into source $R$-blocks.
The number of current background coordinates in each fiber is
finite at fixed $R$. The established classical support
construction, followed by the exact finite-range composite pivot
and ordinary logarithm maps, supplies individually symmetric
strict pieces of
\begin{equation}
 B(a)=A_{M,R}(a)-A_{M,R,0}
 \label{cdc-actual-precision}
\end{equation}
on an independent current-background product ball. Their row
norm has an arbitrary prescribed finite cardinality reserve and
some positive spatial exponent, uniformly in $M$ and volume.
The chain rule includes the classical current contracted against
both chart-curvature terms, and the full absolute gauge Hessian.
These terms are not omitted away from the physical branch.

The actual stationary field is the pre-gauge supported
representative followed by the fixed-$R$ internal contour gauge.
Finite-range Banach-valued analytic composition therefore gives
\eqref{cdc-map-input} for $a_*(w)$ with any fixed finite
cardinality reserve, a positive exponent and $s(r)=O_R(r)$.
Its output fiber is the entire current field in an $R$-block,
with its maximum norm inside the support sum. The lemma gives
a symmetric strict $w$-precision with row norm $O_R(r)$.
The outer radius $b_{\rm cl}$ here is a classical radius; it is
not the inherited quantum radius reused at every step.

Now group source roots and next-background cells into odd
$J$-blocks. For a projected support $\varrho_J Z$, the established
geometric bounds give
\begin{equation}
 \tau_{\rm mac}(\varrho_JZ)
       \leq\frac2J\tau(Z)+d|Z|,\qquad
 \ell_{\rm mac}(\varrho_JZ)
       \leq\frac4J\tau(Z)+2d|Z|.
 \label{cdc-projection}
\end{equation}
The second uses metric spanning-tree length $\ell$, and is
conservative. Grouping matrix rows costs at most $J^d$:
sum the old rows in the macrocell before taking its block norm.
For the target strict Gaussian exponent $\gamma=8$, it is
sufficient in $d=4$ to choose
\begin{equation}
 \frac{32}{J}\leq\kappa',\qquad A'\geq64.
 \label{cdc-precision-reserve}
\end{equation}
This produces the required macro precision norm $O_{R,J}(r)$.
The separately established covariance decay is regrouped at the
same $J$; increasing $J$ if necessary supplies the fixed
Gaussian spatial exponents.

For clarity, the small-radius choice is made \emph{after} $R,J$.
The real baseline retains its spectral bounds
$\gamma_R I\leq A_{M,R,0}\leq B_RI$ under grouping. If
$\nu(r)$ bounds the substituted precision perturbation in the
macro weighted row norm, then
\begin{equation}
 \nu(r)\leq\gamma_R/2
 \quad\Longrightarrow\quad
 \|A(w)^{-1}-A_0^{-1}\|
             \leq2\gamma_R^{-2}\nu(r).
 \label{cdc-covariance-comparison}
\end{equation}
Hence the sufficient additional condition
\begin{equation}
 \nu(r)\leq\frac{\theta\gamma_R^2}{4B_R}
 \label{cdc-relative-neighborhood}
\end{equation}
puts the covariance within relative norm $\theta/2$ of the real
baseline, leaving the other half for the Gaussian Cauchy
perturbation. The strict density-support construction has
additional positive smallness thresholds in this same weighted
precision norm. Since $\nu(r)=O_{R,J}(r)$, all can be met by
choosing a single final $b\leq r$ after the finite geometry and
$\theta$ have been fixed. For example
$\theta=1/(4D_G)$ uses the actual finite macro Gaussian
dimension $D_G$. This choice does not require replacing the
inherited field radius by a smaller fraction after each
stationary substitution.

\subsection{Scalar exponents with spectator fluctuation variables}
\label{cdc-scalar-section}

Let $V_Y(a|_Y,\xi|_Y)$ be a classical or measure exponent family,
holomorphic in the independent background, measurable on the
actual retained cutoff, with
\begin{equation}
 \sup_i\sum_{Y\ni i}e^{\sigma\tau(Y)+A|Y|}
             \|V_Y\|_{b_{\rm cl},{\rm cutoff}}
 \leq\epsilon_{\rm cl}.
 \label{cdc-scalar-input}
\end{equation}
Only $Y$ is read in the fluctuation variables. It need not be
connected. Empty scalar vacuum terms, if present, are tracked
separately. Substitute \eqref{cdc-map-input} in the background
slots while keeping $\xi$ as spectator data.

\begin{lemma}[Joint macro carriers for the base factors]
\label{cdc-scalar-lemma}
Fix a target cardinality exponent $\mathsf B\geq0$ and
$\delta>0$. With $t=s(r)/b_{\rm cl}<1$, suppose
\begin{equation}
 \frac{2(\mathsf B+\delta)}J\leq\min(\sigma,\kappa),\qquad
 a\geq d(\mathsf B+\delta),\qquad
 A\geq d(\mathsf B+\delta)+\alpha_D(t).
 \label{cdc-scalar-budget}
\end{equation}
There is a normally convergent decomposition into connected
strict joint $(w,\xi)$ macro carriers $K$ such that
\begin{equation}
 \sum_YV_Y(a_*(w)|_Y,\xi|_Y)=\sum_KV'_K(w_K,\xi_K),
 \qquad
 \sup_z\sum_{K\ni z}e^{\mathsf B|K|}
                         \|V'_K\|
 \leq e^{\mathsf B}J^d S_\delta\epsilon_{\rm cl},
 \label{cdc-scalar-output}
\end{equation}
where $S_\delta=\sup_z\sum_y e^{-\delta d(z,y)}$.
\end{lemma}
\begin{proof}
Taylor expand in the $D|Y|$ scalar background coordinates only.
Retain $Y$ as a fluctuation carrier even in the zero Taylor
coefficient. Join $\varrho_JY$ by a deterministic Steiner graph;
for each substituted support $X$ rooted at $k\in Y$, attach a
deterministic Steiner graph spanning
$\{\varrho_Jk\}\cup\varrho_JX$. Their union is a connected
joint carrier $K$. If $T$ is the sum of their edge lengths, then
\begin{align}
 |K|&\leq1+T,\nonumber\\
 T&\leq\frac2J\tau(Y)+d|Y|
       +\sum_{\rm occurrences}
          \left[\frac2J\tau(\{k\}\cup X)+d|X|\right].
 \label{cdc-joint-tree}
\end{align}
The carrier includes every background argument and the entire
projected fluctuation carrier, regardless of their overlap.
For $z\in K$, $d(z,\varrho_JY)\leq T$. Thus
\begin{equation}
 e^{\mathsf B|K|}\mathbf1_{z\in K}
 \leq e^{\mathsf B}e^{-\delta d(z,\varrho_JY)}
                                      e^{(\mathsf B+\delta)T}.
 \label{cdc-reroot}
\end{equation}
The spatial and cardinal weights in
\eqref{cdc-scalar-budget} pay the last factor. The scalar
multiindex sum, now including degree zero, is exactly
$(1-t)^{-D|Y|}=e^{\alpha_D(t)|Y|}$, paid by the remaining
input cardinal reserve. Finally use
\begin{equation}
 e^{-\delta d(z,\varrho_JY)}
 \leq\sum_{y\in\varrho_JY}e^{-\delta d(z,y)}.
 \label{cdc-anchor-sum}
\end{equation}
For each projected root $y$ there are at most $J^d$ source
sites. The original anchored norm and the lattice sum
$S_\delta$ yield \eqref{cdc-scalar-output}. These estimates
are positive majorants uniform over the actual real cutoff;
normal convergence gives background holomorphy. All deterministic
support assignments are fixed before choosing the fields.
\end{proof}

The zero background Taylor term can depend on $\xi$ and is then
an interaction, not a vacuum constant. The proof does not infer
a strict family from a bound on the summed scalar functional.
It starts with the actual coefficient bound
\eqref{cdc-scalar-input}.

\subsection{Actual classical and measure inputs and scope}
\label{cdc-matching-section}

At fixed $R$, the already established classical support theorem
applies to the action after the actual composite chart and
ordinary-log conversion. Its Taylor remainder after subtracting
constant, linear and quadratic terms has arbitrarily weighted
strict coefficients. On the physical stationary branch the
subtracted linear term vanishes, so this extension agrees with
the exact centered integral. The full absolute gauge action is
treated the same way; no off-stationary residual is discarded.

The local ambient pivot logarithms, Haar logarithms and gauge
remainder also have actual independent-background coefficient
families. Composite one-pivot incidence keeps their ranges
finite at fixed $R$. The retained link cutoff is unchanged.
With its link threshold denoted by $T$, these families retain
the respective small factors
\begin{equation}
 O_{R,J}\bigl(|g|T^3+|g|T+g^2T^2\bigr)
 \label{cdc-base-smallness}
\end{equation}
after the scalar substitution, provided the fixed chart and
Taylor domains hold. Equal-carrier exponents are grouped before
exponentiation. If their anchored exponent bound is $v$, it
bounds each individual exponent and gives the multiplier norm
bound $e^v v$ for $e^{V'_K}-1$.

Background-only pivot-delta factors are treated by the same
scalar substitution. The Gaussian determinant is expanded in
closed chains using the already substituted strict precision;
its field-independent normalization remains separate. The
classical outer radii and reserves are fixed independently of
the old quantum activity radius.

Thus the precision and base multipliers already read $w$ on
the auxiliary Gaussian graph. Once a separate inherited joint
family meets its budget, the Gaussian expansion outputs
$w$-activities directly on next-current $J$-cells. There is no
second transport through $a_*(w)$ after integration. This changes
only the order of normally convergent support substitution;
the exact integration variables, covariance, cutoff and
normalization are unchanged.

This section is not the integrated quantum-step theorem. It
does not estimate inherited activity amplitudes, prove a
contractive recursion, identify a running coupling or control
large fields. All equality statements start on their common
local analytic domains and extend only on the established
holomorphic continuations.
